\begin{abstract}
We develop and implement certified improvements at two interfaces of the
ProveIt unknot-recognition programme. The principal geometric result is an
exact polynomial reduction of optimal vertex-link-peeled Euler scoring,
over every subbatch of a fixed tetrahedron-disjoint family of coherent
$3$--$2$ Pachner moves, to maximum-weight matching. A minimum-omission
optimum omits at most one move per relevant interior vertex. The maintained
pulling construction has exactly two interior vertices, so its optimum
omits at most two moves and its matching stage has a linear arithmetic-time
selector.
The conflict graph of all distinct legal central edges has degree at most
nine; consequently a packed optimum retains at least
$\max\{0,\lceil M/10\rceil-2\}$ of $M$ current candidates. We also prove
relative-boundary commutation, implement one simultaneous certificate, and
maintain exact peeling minima with deterministic balanced trees. A genuine
32-tetrahedron solid-torus example shows why adding singleton peeled scores
would be incorrect.

On the linear-programming side, a primal--dual pair certifies the entire
active support of a rational nonnegative matching cone. A homogeneous lift
produces that pair or an ordinary small negative certificate. An adaptive
quadrilateral search has at most $2^{\rho+1}-1$ nodes in terms of an intrinsic
active ambiguity rank $\rho$. However, explicit local normal-coordinate
directions prove $\rho\ge t-v$ at unrestricted anchored roots, and
$\rho\ge16\max\{1,n\}-2$ on the maintained $n$-crossing pulling source.
This is a structural obstruction to obtaining a universal polylogarithmic
bound from active support alone.

The additive implementation includes independent consumers, exact replay
fixtures, paired timing data, and a pinned integration patch. Four
64-move workloads show approximately $117$--$123$-fold producer improvements
over literal sequential replay, while the enlarged support LP regresses
against small unary probes. These are subroutine results. We identify the
remaining witness-coverage, transition-count, anchor, and bit-complexity
theorems required for a complete quasipolynomial recognizer, and propose a
concrete research programme to address them.
\end{abstract}

